\documentclass[a4paper,10pt]{report}\input{../0.Préambule/Préambule_cours_prof}\begin{document}

\pagestyle{empty}
\newpage
\noindent NOM, Prénom et classe : \dotfill

\section*{Statistiques \hspace{3cm} Sujet A}
\addcontentsline{toc}{section}{Statistiques Sujet A}

\vspace{-2mm}
\noindent \textit{L'usage de la calculatrice est \textbf{autorisée}.} \hfill \textit{Durée : 55 minutes}

\paragraph{Exercice 1} : \quad Donner la définition de la fréquence en statistique. \hfill \textbf{2 points} 

\paragraph{Exercice 2} : \hfill \textbf{5 points}

Dans les classes de 5$^e$ TANDEM et de 5$^e$ COUBERTIN, $32$ élèves sont demi-pensionnaires dont $14$ sont en 5$^e$ COUBERTIN. Les $12$ autres élèves de la classe sont externes comme $6$ élèves de la classe de 5$^e$ TANDEM.

\medskip
\begin{enumerate}
\item Complète le tableau suivant.
\item Combien d'élèves y a-t-il en 5$^e$ TANDEM.
\item Quelle la fréquence d'élève demi-pensionnaire sur ces deux classes.
\item Quelle est la fréquence d'élève externe en 5$^e$ COUBERTIN.
\end{enumerate}

\vspace{-3cm}
\begin{flushright}
\arrayrulecolor{black}
\renewcommand{\arraystretch}{1.5}
\begin{tabularx}{6cm}{|p{1.5cm}|*{3}{>{\centering \arraybackslash}X|}}
\cline{2-4}
\multicolumn{1}{c|}{} & 5$^e$ T & 5$^e$ C & Total \\\hline
Externes & & &\\\hline
D.P. & & &\\\hline
Total & & &\\\hline
\end{tabularx}
\end{flushright}

\paragraph{Exercice 3} : \hfill \textbf{5 points}

\medskip Bastien a relevé les durées des morceaux de sa compilation de rock symphonique préférée en secondes.

\begin{center}
\noindent $248$ \quad ; \quad $199$ \quad ; \quad $287$ \quad ; \quad $226$ \quad ; \quad $195$ \quad ; \quad $199$ \quad ; \quad $238$ \quad ; \quad $230$ \quad ; \quad $204$ \quad ; \quad $235$ \quad ; \quad $196$ \quad ; \quad $204$ \quad ; \quad $187$ 
\end{center}

\begin{enumerate}
\item Calcule la durée moyenne des morceaux.
\item Détermine la durée médiane.
\end{enumerate}
\noindent Bastien rajoute une musique qui dure $3$ minutes et $25$ secondes. 
\begin{enumerate}
\setcounter{enumi}{2}
\item Que fait la durée moyenne des morceaux ? Explique ton raisonnement.
\end{enumerate}

\paragraph{Exercice 4} : 

\subparagraph{Partie A} : Étude avec un logiciel\hfill \textbf{4 points} 

\medskip
On liste les couleurs des yeux d'une population d'élève de cinquième. On dresse dans un logiciel le tableau ci contre avec les résultats obtenus :

\noindent \begin{minipage}{8.5cm}
\begin{enumerate}
\item Donner le nom d'un tel logiciel.
\item Quelle est le contenu de la case $D1$.
\item Que contient la ligne $2$.
\item Quelle formule pourrait-on écrire dans la case $E2$ ?
\end{enumerate}
\end{minipage}
\hspace{10mm}
\begin{minipage}{7cm}
\begin{tableur}{5}
& Couleur & Bleu & Vert & Marron & Total \\\hline
& Effectif & 24 & 24 & 35 & \\\hline
& & & & & \\\hline
\end{tableur}
\end{minipage}

\subparagraph{Partie B} : Étude avec un diagramme en barre \hfill \textbf{4 points} 

\medskip
Sur les même personne, on s'intéresse maintenant à leur jour de naissance. $18$ sont nés un lundi, $5$ un mardi, $4$ un mercredi, $12$ un jeudi, $14$ un vendredi, $8$ un samedi et le reste un dimanche.

\begin{enumerate}
\setcounter{enumi}{4}
\item Dresse un diagramme en bâton permettant de représenter cette série statistiques. 

\textit{Indication :} On prendra $1$ carreau pour $2$ personnes en ordonnées .

On n'oubliera pas de légender les axes et de mettre un titre au graphique.

\item Quelle est la fréquence de personne née un mercredi ?
\end{enumerate}

\newpage
\noindent NOM, Prénom et classe : \dotfill

\section*{Statistiques \hspace{3cm} Sujet B}
\addcontentsline{toc}{section}{Statistiques Sujet B}

\vspace{-2mm}
\noindent \textit{L'usage de la calculatrice est \textbf{autorisée}.} \hfill \textit{Durée : 55 minutes}

\paragraph{Exercice 1} : \quad Donner la définition de la fréquence en statistique. \hfill \textbf{2 points} 

\paragraph{Exercice 2} : \hfill \textbf{5 points}

Dans les classes de 5$^e$ TANDEM et de 5$^e$ COUBERTIN, $32$ élèves sont demi-pensionnaires dont $14$ sont en 5$^e$ COUBERTIN. Les $10$ autres élèves de la classe sont externes comme $11$ élèves de la classe de 5$^e$ TANDEM.

\medskip
\begin{enumerate}
\item Complète le tableau suivant.
\item Combien d'élèves y a-t-il en 5$^e$ TANDEM.
\item Quelle la fréquence d'élève demi-pensionnaire sur ces deux classes.
\item Quelle est la fréquence d'élève externe en 5$^e$ COUBERTIN.
\end{enumerate}

\vspace{-3cm}
\begin{flushright}
\arrayrulecolor{black}
\renewcommand{\arraystretch}{1.5}
\begin{tabularx}{6cm}{|p{1.5cm}|*{3}{>{\centering \arraybackslash}X|}}
\cline{2-4}
\multicolumn{1}{c|}{} & 5$^e$ T & 5$^e$ C & Total \\\hline
Externes & & &\\\hline
D.P. & & &\\\hline
Total & & &\\\hline
\end{tabularx}
\end{flushright}

\paragraph{Exercice 3} : \hfill \textbf{5 points}

\medskip Bastien a relevé les durées des morceaux de sa compilation de rock symphonique préférée en secondes.

\begin{center}
\noindent $278$ \quad ; \quad $199$ \quad ; \quad $287$ \quad ; \quad $226$ \quad ; \quad $195$ \quad ; \quad $199$ \quad ; \quad $278$ \quad ; \quad $230$ \quad ; \quad $234$ \quad ; \quad $255$ \quad ; \quad $196$ \quad ; \quad $224$ \quad ; \quad $197$ 
\end{center}

\begin{enumerate}
\item Calcule la durée moyenne des morceaux.
\item Détermine la durée médiane.
\end{enumerate}
\noindent Bastien rajoute une musique qui dure $4$ minutes et $25$ secondes. 
\begin{enumerate}
\setcounter{enumi}{2}
\item Que fait la durée moyenne des morceaux ? Explique ton raisonnement.
\end{enumerate}

\paragraph{Exercice 4} : 

\subparagraph{Partie A} : Étude avec un logiciel\hfill \textbf{4 points} 

\medskip
On liste les couleurs des yeux d'une population d'élève de cinquième. On dresse dans un logiciel le tableau ci contre avec les résultats obtenus :

\noindent \begin{minipage}{8.5cm}
\begin{enumerate}
\item Donner le nom d'un tel logiciel.
\item Quelle est le contenu de la case $C1$.
\item Que contient la ligne $2$.
\item Quelle formule pourrait-on écrire dans la case $E2$ ?
\end{enumerate}
\end{minipage}
\hspace{10mm}
\begin{minipage}{7cm}
\begin{tableur}{5}
& Couleur & Bleu & Vert & Marron & Total \\\hline
& Effectif & 27 & 22 & 35 & \\\hline
& & & & & \\\hline
\end{tableur}
\end{minipage}

\subparagraph{Partie B} : Étude avec un diagramme en barre \hfill \textbf{4 points} 

\medskip
Sur les même personne, on s'intéresse maintenant à leur jour de naissance. $18$ sont nés un lundi, $5$ un mardi, $7$ un mercredi, $12$ un jeudi, $11$ un vendredi, $8$ un samedi et le reste un dimanche.

\begin{enumerate}
\setcounter{enumi}{4}
\item Dresse un diagramme en bâton permettant de représenter cette série statistiques. 

\textit{Indication :} On prendra $1$ carreau pour $2$ personnes en ordonnées .

On n'oubliera pas de légender les axes et de mettre un titre au graphique.

\item Quelle est la fréquence de personne née un mercredi ?
\end{enumerate}
\end{document}